A large language model was used in the preparation of this work, and we describe
its role precisely.

\emph{Where it was used.} (i)~Writing and editing the manuscript, including
restructuring, translation, and prose revision. (ii)~Writing the analysis code in
\texttt{scripts/paper/} --- the statistics module, the artifact loaders, and the
table generator. (iii)~Auditing the committed experiment artifacts, where it
surfaced the evaluator-fallback contamination, the degenerate recovery columns,
the 800-record matrix being 400 pairs, and the architecture table's pooled
fault bucket.

\emph{Where it was not used.} The fault-injection harness, the agent
implementation, and the experiment runs themselves predate this assistance. The
model did not design the experiments, select the tasks, or produce any
experimental result. It did not fabricate, interpolate, or adjust any number: every
quantity in the paper is computed by \texttt{scripts/paper\_tables.py} from the
committed artifacts, and that script is released.

\emph{Verification.} Two safeguards apply to the analysis code, which is the part
where an error would silently change a result. First, the statistical routines
carry unit tests with hand-computed expected values, and one such test caught a
real defect during preparation --- an incorrect search bracket in the minimum
detectable effect calculation that made the function return an undefined value for
every input. Second, the main tables are checked against independently committed
summaries: the pooled paired success rates reproduce
\texttt{experiments/analysis/official\_subset\_v1\_final/official\_summary.csv}
cell for cell. The author reviewed and verified all generated output and takes
responsibility for the content.
